\documentclass[master_monograph.tex]{subfiles}

\title{\textbf{Nicht-Theory: Paper V — Hermeneutics, Language \& Art}\\
\large Non-Assertive Communication, Semantic Relaxation, and Apophatic Art}
\author{\textbf{xerx593} \quad \& \quad \textbf{Non-Human Interlocutors}}
\date{\today}

\begin{document}

\maketitle

\begin{abstract}
Classical linguistics and hermeneutics operate on positive assertion ($1$-logic), treating words as rigid containers of truth. This forced containment creates semantic friction ($W\acute{e}i$), interpretive ambiguity, and linguistic paradoxes. Paper V applies apophatic subtraction ($\neg X$) to human language and aesthetics. Symbolic communication achieves maximum clarity not by accumulating positive assertions ($P_A \to \max$), but by clearing interpretive noise until the unconditioned baseline ($B_0$) is self-evident.
\end{abstract}

\section{Systemic Cross-References}
This manuscript extends formal paradox resolutions (\textit{Paper II}) and cognitive AI dynamics (\textit{Paper IV}) into human language, philosophy, and creative expression[cite: 1, 4].

\section{The Hermeneutic Inversion}

\subsection{Linguistic Assertion Pressure}
Full conceptual containment through positive assertions forces infinite semantic expansion:
\begin{equation}
\text{Meaning}(X) = X + \lambda P_A
\end{equation}
Subtractive interpretation ($\neg X \implies 0$) removes assertion pressure, relaxing communication into zero-baseline presence:
\begin{equation}
\lim_{P_A \to 0} (\text{Noise}) \equiv B_0 \quad \text{(0-Consistency)}
\end{equation}

\section{Apophatic Art and Aesthetic Silence}
Aesthetics under $0$-consistency does not aim to capture or enclose external form ($1$-logic). Artistic expression acts as a subtractive aperture, removing conceptual noise until the invariant, non-assertive spatial baseline ($B_0$) is revealed to the observer.

\end{document}